% Nouveau EGA V (cote 25): archive cote 25, site page 133 (PDF page 134).
% typescript. Draft compared: none.
% AI-transcribed by Claude Opus 5.5 (claude-opus-5-5, adaptive thinking, effort high), Message Batches API, batch ega25b (msgbatch_01A7n5YuXubUPNK1zsepi8RG), 2026-11-13; original AI transcription, subsequently checked against both supplied scan parts; see PAGES_CHECKED.tsv and CHECK_DECISIONS.tsv.
\scanpage{25/133}
\mspage{\hand{36}}

un sous-k-espace vectoriel\var{\inserted[typed]{E}}{ E }{insertion}{ega25_25_133-v0001}de l'espace des \var{\deleted{\var{\inserted{pseudo}}{pseudo}{insertion}{ega25_25_133-v0003}}}{}{deletion}{ega25_25_133-v0002} fonctions \var{\deleted{rationelles}}{}{deletion}{ega25_25_133-v0004} \var{\inserted{méromorphes}}{méromorphes}{insertion}{ega25_25_133-v0005} \var{\deleted{\var{\inserted[typed]{strictes}}{strictes}{insertion}{ega25_25_133-v0007}}}{}{deletion}{ega25_25_133-v0006} \var{\xout{xxxxxxx}}{}{typewriter strike-out}{ega25_25_133-v0008}
sur X, tel que pour tout $\hand{\phi} \hand{\in} E-(0)$, $\hand{\phi}$ soit \var{\inserted{régulier}}{régulier}{insertion}{ega25_25_133-v0009} \var{\deleted{inversible dans l'\var{\xout{xxxxxxx}}{}{typewriter strike-out}{ega25_25_133-v0011}}}{}{deletion}{ega25_25_133-v0010}
\var{\deleted{anneau des fonctions rationelles}}{}{deletion}{ega25_25_133-v0012} \var{\deleted{\var{\inserted[typed]{strictes}}{strictes}{insertion}{ega25_25_133-v0014}}}{}{deletion}{ega25_25_133-v0013} \var{\xout{xxxxxx}}{}{typewriter strike-out}{ega25_25_133-v0015}, i.e. div($\hand{\phi}$) est défini \var{\deleted{comme}}{}{deletion}{ega25_25_133-v0016}
\var{\deleted{diviseur de Cartier,}}{}{deletion}{ega25_25_133-v0017}\hand{\unsure{,|)}} et que $\hand{\triangle}$ soit l'ensemble des div($\hand{\phi}$) pour $\hand{\phi} \hand{\in} E-(0)$.
c) Soient E, E' deux sous-\var{\inserted[typed]{k-}}{k-}{insertion}{ega25_25_133-v0018}espaces vectoriels de l'espace des
fonctions rationelles \var{\inserted[typed]{strictes}}{strictes}{insertion}{ega25_25_133-v0019} \var{\xout{xxxxxxxx}}{}{typewriter strike-out}{ega25_25_133-v0020} sur X, \var{\xout{xxxxx}}{}{typewriter strike-out}{ega25_25_133-v0021} satisfaisant la condition
envisagée dans b), alors les ensembles de diviseurs $\hand{\triangle}$, $\hand{\triangle}'$ qu'ils
définissent respectivement sont égaux sss il existe une \var{\inserted{pseudo-}}{pseudo-}{insertion}{ega25_25_133-v0022} fonction
\var{\deleted{rationelle stricte}}{}{deletion}{ega25_25_133-v0023}e \var{\deleted{inversible}}{}{deletion}{ega25_25_133-v0024} \var{\inserted{régulière}}{régulière}{insertion}{ega25_25_133-v0025} $\hand{\phi}$ sur X, telle que $E' = \hand{\phi} E$ . Si $E\hand{\neq}(0)$\var{\inserted{(i.e. $\triangle \neq \emptyset$)}}{(i.e. $\triangle \neq \emptyset$)}{insertion}{ega25_25_133-v0026}, une
telle $\hand{\phi}$ est déterminée modulo multiplication par un élément c de $k^{\hand{*}}$.

La démonstration est un exercice facile, enutilisant \hand{4}.1. , et je
me dispense d'en écrire la démonstration, sauf réclamation de ta part. Il
\var{\margin{||}}{\placedmargin{||}}{marginal note placed}{ega25_25_133-v0027} me sembled'ailleurs que \hand{4}.4. viendrait avantageusement dès avant \hand{4}.3.,
étanttechniquement nettement plus trivial. Noter d'ailleurs que lorsque
X est géom intègre, la condition sur E énoncée dans b) devient \uline{vide}.
La restriction mise au début de b) tient au fait qu'il se pourrait autre-%
ment que la condition énoncée pour b) ne soit plus vérifiée après passage
à la clôture algébrique de k (il est facile d'en donner des exemples\var{\xout{xx}}{}{typewriter strike-out}{ega25_25_133-v0028},
en toute caractéristique, et même lorsque k est séparablement clos en
c\var{\xout{\unsure{ar}}}{}{typewriter strike-out}{ega25_25_133-v0029} p \hand{$>$} 0). Pour bien faire, on devrait énoncer b) sans condition supplé-%
mentaire sur k ou X, mais en énoncant la condition sur E en passant à la
clôture algébrique de k (et signaler que si X est géom intègre, \unsure{oette|cette}
condition devient vide). \hand{-- Par abus de langage, un}
\hand{\uline{ens.} $\triangle$ de diviseurs de la forme \unsure{Ens|En}(D) sera}
\hand{encore parfois appelé système linéaire de diviseurs}
\hand{sur X ;}

\mspage{\hand{133}}
